\section{Results}\label{sec:results}

\subsection{Gate level}\label{sec:res_gate}
\Cref{tab:gates} summarises the gates over \NTransitions{} transitions each. The original gate detects almost every motion transition (recall \GateRecallFD) but triggers on \GateFprFD{} of the transitions without motion, almost entirely because of camera jitter, where it fires on \GateFprJitterFD{} of them. \FDS{} keeps the recall (\GateRecallFDS), reduces the false-positive rate to \GateFprFDS{} (95\% interval \GateFprCIFDS) and raises the precision from \GatePrecFD{} to \GatePrecFDS. MOG2 misses more motion (recall \GateRecallMOG) because it needs frames to build its model and partly absorbs slow motion; it adapts to illumination steps only gradually, and it is the most expensive gate (\GateMsMOG\,ms per frame).

\begin{table}[t]
\centering
\caption{Gate-level results on the controlled benchmark, pooled over \NSeeds{} realisations (\NTransitions{} transitions per gate). Recall and FPR are fractions of the transitions with and without motion; the per-scenario columns single out the hardest cases. Costs per $640\times640$ frame on one CPU thread.}
\label{tab:gates}
\resizebox{\linewidth}{!}{\begin{tabular}{@{}lccccccc@{}}
\toprule
Gate & Recall & FPR & Precision & Recall & FPR & FPR & Cost \\
 & (all) & (all) & (all) & distant walker & illumination & jitter & (ms/frame) \\
\midrule
FD (original) & 0.996 & 0.206 & 0.832 & 1.000 & 0.022 & 0.996 & 1.43 \\
FD-S (stabilised) & 0.996 & 0.014 & 0.987 & 1.000 & 0.000 & 0.069 & 7.25 \\
MOG2 & 0.944 & 0.114 & 0.895 & 0.955 & 0.236 & 0.226 & 10.51 \\
\bottomrule
\end{tabular}
}
\end{table}

\begin{figure}[t]
\centering
\includegraphics[width=\linewidth]{fig_envelope}
\caption{Operating envelope of \FD{} and \FDS{} (3 realisations $\times$ 30 transitions per point; shaded: 95\% Wilson intervals). (a)~Noise; bands mark the predicted onset for per-pixel exceedance probabilities between 0.01 and 0.1, dotted lines the percolation prediction. (b)~Illumination steps; the dashed line marks $\Delta=\tau$. (c)~Camera translation between consecutive frames. (d)~Walking persons of three heights; dotted lines mark the area-model bound $v^*$ where it is positive.}
\label{fig:envelope}
\end{figure}

The sweeps in \cref{fig:envelope} test the predictions of \cref{sec:envelope}:
\begin{itemize}
\item \emph{Noise.} \FD{} switches from never to always active at $\sigma=\MeasSigmaFD$ (50\% crossing), where the percolation criterion predicts $\PredSigmaPcFD$. \FDS{} switches at $\sigma=\MeasSigmaFDS$, predicted $\PredSigmaPcFDS$. The rectification bias thus accounts quantitatively for the difference in noise tolerance.
\item \emph{Illumination.} \FD{} fires from $\Delta=\MeasIllumFD$ on, which matches the prediction $\Delta>\tau=20$ up to the integer rounding of the OpenCV pipeline. \FDS{} tolerates steps of up to 30 grey levels and fires beyond (crossing at \MeasIllumFDS). For large steps, the neighbourhoods of saturated regions change by less than the offset, and after smoothing such partially clipped pixels escape the saturation mask. The illumination scenario of the benchmark (steps of $+32$ and $-50$) does not trigger \FDS, but its immunity is not complete.
\item \emph{Camera jitter.} \FD{} fires from translations of \MeasJitterFD\,px on, the same order of magnitude as the estimate $\delta^*\approx\PredJitter$\,px. \FDS{} suppresses translations of 1.5\,px and more completely but fires in a band around 0.5--1\,px. Such shifts already exceed the threshold at strong edges, yet they are too small for a reliable half-resolution estimate, so the compensation is rejected or inaccurate. This band explains the residual false-positive rate of \GateFprJitterFDS{} in the jitter scenario.
\item \emph{Size and speed.} All three persons trigger \FD{} already at 1\,px/frame, far below the area-model bounds of \PredSpeedA{} and \PredSpeedB\,px/frame for heights of \HeightA{} and \HeightB\,px. The textured interior of the silhouette changes many more pixels than the two edge strips, so the bound is tight only for weakly textured objects. \FDS{} is slightly less sensitive to the smallest persons at 1\,px/frame, because its signed difference lacks the positive bias that helps \FD.
\end{itemize}

\noindent \Cref{fig:pipeline} illustrates the difference on a frame of the jitter-and-walk scenario: \FD{} marks edges all over the scene, whereas the change mask of \FDS{} is dominated by the walking person.

\begin{figure}[t]
\centering
\includegraphics[width=\linewidth]{fig_pipeline}
\caption{A frame of the jitter-and-walk scenario (estimated camera shift \PipelineShift\,px). \FD{} marks edges all over the scene; \FDS{} compensates the shift and its mask is dominated by the walking person, with small residuals at the strongest edges. In (d) the red box is the YOLOv8m detection and the green boxes are the \FDS{} change regions.}
\label{fig:pipeline}
\end{figure}

\subsection{System level}\label{sec:res_system}
\begin{table}[t]
\centering
\caption{System-level results on the controlled benchmark ($K=\RefreshK$, $M=1$). Stale frames: frames without a person on which a box is reported. Latencies in frames relative to the every-frame detector. Time: measured gate costs plus measured detector costs of the invoked frames.}
\label{tab:policies}
\resizebox{\linewidth}{!}{\begin{tabular}{@{}lccccccccc@{}}
\toprule
Policy & $r$ & Recall & Stale & Mean & Max & Entry & Exit & Time & Speed- \\
 & & & frames & IoU & age & lat. & lat. & (s) & up \\
\midrule
Every frame & 1.000 & 0.998 & 0 & 0.981 & 0 & 0 & 0 & 628.3 & 1.00 \\
Fixed rate, $N=2$ & 0.500 & 0.997 & 0 & 0.938 & 1 & 0 & 0 & 313.3 & 2.01 \\
Gated, FD (original) & 0.629 & 0.997 & 0 & 0.981 & 14 & 1 & 0 & 395.0 & 1.59 \\
Gated, FD-S (stabilised) & 0.542 & 0.997 & 0 & 0.981 & 14 & 1 & 0 & 342.6 & 1.83 \\
Gated, MOG2 & 0.557 & 0.997 & 5 & 0.973 & 14 & 1 & 6 & 356.8 & 1.76 \\
\bottomrule
\end{tabular}
}
\end{table}

\Cref{tab:policies} compares the policies at the default setting. With \FDS{} the detector runs on \PolRatioFDS{} of the frames (\PolCallsFDS{} calls), and the system still matches the every-frame output: recall \PolRecallFDS{} vs.\ \PolRecallEvery, the same mean IoU (\PolIouFDS), no stale frames and no exit latency. The maximum age is $K-1=14$ frames, as guaranteed by Proposition~\ref{prop:guarantees}. The one-frame entry latency arises because the first visible sliver of an entering person changes fewer than $A_{\min}$ pixels. \FD{} reaches the same accuracy but needs \PolCallsFD{} calls because of its jitter triggers (speed-up \PolSpeedFD{} instead of \PolSpeedFDS). MOG2 did not register the person leaving the image, so its last box remained for \PolExitMOG{} frames until the refresh.

\Cref{fig:scenarios} and \cref{tab:scenario_results} break the invocation ratio down by scenario. Static scenarios cost $\lceil T/K\rceil/T=0.067$ with every gate, the lower bound of Proposition~\ref{prop:guarantees}, except for \FD{} under camera jitter and MOG2 under illumination steps. Scenarios with continuous motion reach the upper bound $1/M=1$.

\begin{figure}[t]
\centering
\includegraphics[width=\linewidth]{fig_scenarios}
\caption{Invocation ratio per scenario ($K=\RefreshK$, $M=1$). The dotted line marks the lower bound $1/K$, the dashed line the upper bound $1/M$.}
\label{fig:scenarios}
\end{figure}

\begin{table}[t]
\centering\small
\caption{Invocation ratio and frame recall per scenario ($K=\RefreshK$, $M=1$); recall is undefined without a person.}
\label{tab:scenario_results}
\begin{tabular}{@{}lcccccc@{}}
\toprule
 & \multicolumn{3}{c}{Invocation ratio $r$ ($K=15$)} & \multicolumn{3}{c}{Frame recall} \\
\cmidrule(lr){2-4}\cmidrule(l){5-7}
Scenario & FD & FD-S & MOG2 & FD & FD-S & MOG2 \\
\midrule
Walk & 1.000 & 1.000 & 0.956 & 1.000 & 1.000 & 1.000 \\
Slow walk & 1.000 & 1.000 & 0.933 & 1.000 & 1.000 & 1.000 \\
Stop and go & 0.689 & 0.689 & 0.700 & 1.000 & 1.000 & 1.000 \\
Enter, stop, exit & 0.411 & 0.411 & 0.433 & 0.973 & 0.973 & 0.973 \\
Stationary person & 0.067 & 0.067 & 0.067 & 1.000 & 1.000 & 1.000 \\
Distant walker & 1.000 & 1.000 & 0.956 & 1.000 & 1.000 & 1.000 \\
Empty scene & 0.067 & 0.067 & 0.067 & -- & -- & -- \\
Illumination steps & 0.067 & 0.067 & 0.278 & -- & -- & -- \\
Camera jitter & 0.989 & 0.122 & 0.222 & -- & -- & -- \\
Jitter + walk & 1.000 & 1.000 & 0.956 & 1.000 & 1.000 & 1.000 \\
\bottomrule
\end{tabular}

\end{table}

\Cref{fig:tradeoff}(a,b) compares the policy families across budgets. At IoU $\ge0.5$, fixed-rate subsampling is as good as gating with $M=1$ (\FixedRecallAtFDS{} at the budget of \FDS), but it degrades localisation (mean IoU \FixedIouAtFDS{} vs.\ \PolIouFDS). With larger $M$ the gated policies dominate: at $M=2$ \FDS{} keeps a recall of \PolRecallFDSMTwo{} at $r=\PolRatioFDSMTwo$ (fixed rate: \FixedRecallAtFDSMTwo), and at $M=4$ it keeps \PolRecallFDSMFour{} at $r=\PolRatioFDSMFour$ (fixed rate: \FixedRecallAtFDSMFour).

\begin{figure}[t]
\centering
\includegraphics[width=\linewidth]{fig_tradeoff}
\caption{Accuracy versus detector budget. (a,b)~Controlled benchmark: gated policies with $K=\RefreshK$ and $M\in\{1,2,3,4,6,8\}$ against fixed-rate subsampling; (b)~shows the mean IoU of matched frames, i.e.\ localisation quality. (c)~Real video with $K=10$ and $M\in\{1,\dots,6\}$: the scene is almost always in motion, and all gated policies coincide with fixed-rate subsampling.}
\label{fig:tradeoff}
\end{figure}

\subsection{Cost model and online behaviour}\label{sec:res_cost}
On \EtoeFrames{} frames of the benchmark, the online system with \FD{} made \EtoeCallsFD{} detector calls and took \EtoeMeasuredFD\,s, against \EtoeEstimatedFD\,s estimated by the cost model (\EtoeErrFD\%). With \FDS{} it made \EtoeCallsFDS{} calls in \EtoeMeasuredFDS\,s against \EtoeEstimatedFDS\,s (\EtoeErrFDS\%). The decisions were identical to the offline replay (\EtoeMismatchFD{} and \EtoeMismatchFDS{} mismatches). The estimate is slightly pessimistic because the per-call costs were recorded while the machine was more heavily loaded (mean \YoloMs\,ms per call).

\subsection{Real video}\label{sec:res_video}
On \code{vtest.avi} every frame contains persons (\VtPersons{} on average), and the gates are active on \VtActiveFD{} (\FD), \VtActiveFDS{} (\FDS) and \VtActiveMOG{} (MOG2) of the frames (\cref{tab:video}). With $M=1$ the invocation ratio is therefore close to one and gating saves nothing; the gate's own cost even makes it marginally slower (speed-up \VtSpeedFDS{} with \FDS). With $M=2$ and $M=3$ the gated policies coincide with fixed-rate subsampling (recall \VtRecallFDSTwo{} at $r=\VtRatioFDSTwo$ and \VtRecallFDSThree{} at $r=\VtRatioFDSThree$), exactly as Proposition~\ref{prop:guarantees} predicts for continuous motion. Agreement at a given budget is lower than on the controlled benchmark: at \VtFps\,FPS a skipped frame is already 100\,ms old, and the detector flickers on small, partly occluded pedestrians.

\begin{table}[t]
\centering\small
\caption{Real video (\code{vtest.avi}, \VtFrames{} frames): agreement with YOLOv8m run on every frame ($K=10$).}
\label{tab:video}
\begin{tabular}{@{}lccccc@{}}
\toprule
Policy & $r$ & Recall & Precision & F1 & Speed-up \\
\midrule
Every frame & 1.000 & 1.000 & 1.000 & 1.000 & 1.00 \\
Fixed rate, $N=2$ & 0.501 & 0.952 & 0.954 & 0.953 & 2.00 \\
Fixed rate, $N=3$ & 0.333 & 0.855 & 0.854 & 0.854 & 3.00 \\
Fixed rate, $N=4$ & 0.250 & 0.731 & 0.734 & 0.732 & 4.00 \\
Gated, FD (original), $M=1$ & 0.997 & 1.000 & 1.000 & 1.000 & 1.00 \\
Gated, FD (original), $M=2$ & 0.501 & 0.952 & 0.954 & 0.953 & 1.99 \\
Gated, FD (original), $M=3$ & 0.333 & 0.855 & 0.854 & 0.854 & 2.98 \\
Gated, FD-S (stabilised), $M=1$ & 0.994 & 1.000 & 1.000 & 1.000 & 0.99 \\
Gated, FD-S (stabilised), $M=2$ & 0.498 & 0.952 & 0.954 & 0.953 & 1.95 \\
Gated, FD-S (stabilised), $M=3$ & 0.333 & 0.853 & 0.853 & 0.853 & 2.88 \\
Gated, MOG2, $M=1$ & 0.999 & 1.000 & 1.000 & 1.000 & 0.98 \\
Gated, MOG2, $M=2$ & 0.501 & 0.952 & 0.954 & 0.953 & 1.91 \\
Gated, MOG2, $M=3$ & 0.333 & 0.855 & 0.854 & 0.854 & 2.82 \\
\bottomrule
\end{tabular}

\end{table}
